\HMTNarrativeHeading{La profundidad de un número}

La generación coinductiva organiza la determinación de un número mediante
prolongaciones compatibles. Una etapa produce un prefijo y conserva las
condiciones para continuarlo. La profundidad puede aumentar manteniendo
la concordancia de todas las lecturas anteriores. El valor aparece al
realizar esa compatibilidad; la historia conserva, además, el modo en que
la precisión fue obtenida.

Las regiones de las coordenadas circular, áurea y exponencial proceden del
mismo desarrollo APP--TRIT--TPK. Sus lectores cumplen funciones
diferenciadas y comparten la estructura que permite relacionarlos. La
extensión triádica de Euler explicita las realizaciones de los regímenes;
el calendario areal--volumétrico determina la autoescala; las emisiones
regionales y sus fronteras conservan las condiciones de evaluación.
La sucesión de estos pasos hace explícito el orden de construcción de las
coordenadas que intervienen después en alfa.

El calendario de vacancias permite comprender una particularidad del
refinamiento. La historia avanza en cada transición, mientras la capacidad
de una lectura puede permanecer fija. La palabra que registra esas
transiciones distingue la capacidad adquirida de la vacancia conservada.
Su balance cuenta ambas contribuciones. La regularidad del calendario y
la aperiodicidad de su itinerario describen aspectos relacionados de una
misma evolución.

La monodromía prolonga esta distinción. Volver a una fase significa
reconocer una posición del ciclo; recuperar una historia significa
restituir también sus vueltas, orientación y registros. El espejo regional
conserva el eje y el agregado de las dos coordenadas que intercambia,
mientras invierte su diferencia. Este dato permite reconstruir el
reparto. La figura nonádica sitúa las fases y sus operaciones; sus marcas
son posiciones del calendario, con una semántica diferente del valor real
de una constante.

La quinta frontera, las regiones posteriores y el retorno se explican,
por tanto, mediante sus acciones sobre el estado. Las figuras regionales
muestran simultáneamente coincidencia de una lectura y diferenciación de
sus antecedentes. Esa diferenciación transmite a la realización
excepcional la información incidencial que utiliza. La generación
numérica continúa así la doble proyección planteada al comienzo: cada
lectura obtiene su valor de una estructura cuya organización permanece
disponible para las otras realizaciones.

